\documentclass[10pt,a4paper]{amsart}
\usepackage{amsmath,amssymb,amsthm,graphicx}
% Retained minimal preamble for the Stacks Project chapter flat.tex (More on Flatness).
% Verbatim excerpt of the chapter preamble: Xy-pic setup (native preamble lines 14-18) and the native
% theorem declarations with the mathematical macros (native preamble lines 154-217). Copied mechanically;
% nothing added, removed or redefined. The wrapper class and the amsmath/amssymb/amsthm/graphicx base
% packages are supplied by the assembler, so the class/externaldocument part of the native preamble is omitted.
% For commutative diagrams we use Xy-pic
\usepackage[all]{xy}

% We use 2cell for 2-commutative diagrams.
\xyoption{2cell}
% Theorem environments.
%
\theoremstyle{plain}
\newtheorem{theorem}[subsection]{Theorem}
\newtheorem{proposition}[subsection]{Proposition}
\newtheorem{lemma}[subsection]{Lemma}

\theoremstyle{definition}
\newtheorem{definition}[subsection]{Definition}
\newtheorem{example}[subsection]{Example}
\newtheorem{exercise}[subsection]{Exercise}
\newtheorem{situation}[subsection]{Situation}

\theoremstyle{remark}
\newtheorem{remark}[subsection]{Remark}
\newtheorem{remarks}[subsection]{Remarks}

\numberwithin{equation}{subsection}

% Macros
%
\def\lim{\mathop{\mathrm{lim}}\nolimits}
\def\colim{\mathop{\mathrm{colim}}\nolimits}
\def\Spec{\mathop{\mathrm{Spec}}}
\def\Hom{\mathop{\mathrm{Hom}}\nolimits}
\def\Ext{\mathop{\mathrm{Ext}}\nolimits}
\def\SheafHom{\mathop{\mathcal{H}\!\mathit{om}}\nolimits}
\def\SheafExt{\mathop{\mathcal{E}\!\mathit{xt}}\nolimits}
\def\Sch{\mathit{Sch}}
\def\Mor{\mathop{\mathrm{Mor}}\nolimits}
\def\Ob{\mathop{\mathrm{Ob}}\nolimits}
\def\Sh{\mathop{\mathit{Sh}}\nolimits}
\def\NL{\mathop{N\!L}\nolimits}
\def\CH{\mathop{\mathrm{CH}}\nolimits}
\def\proetale{{pro\text{-}\acute{e}tale}}
\def\etale{{\acute{e}tale}}
\def\QCoh{\mathit{QCoh}}
\def\Ker{\mathop{\mathrm{Ker}}}
\def\Im{\mathop{\mathrm{Im}}}
\def\Coker{\mathop{\mathrm{Coker}}}
\def\Coim{\mathop{\mathrm{Coim}}}

% Boxtimes
%
\DeclareMathSymbol{\boxtimes}{\mathbin}{AMSa}{"02}

%
% Macros for moduli stacks/spaces
%
\def\QCohstack{\mathcal{QC}\!\mathit{oh}}
\def\Cohstack{\mathcal{C}\!\mathit{oh}}
\def\Spacesstack{\mathcal{S}\!\mathit{paces}}
\def\Quotfunctor{\mathrm{Quot}}
\def\Hilbfunctor{\mathrm{Hilb}}
\def\Curvesstack{\mathcal{C}\!\mathit{urves}}
\def\Polarizedstack{\mathcal{P}\!\mathit{olarized}}
\def\Complexesstack{\mathcal{C}\!\mathit{omplexes}}
% \Pic is the operator that assigns to X its picard group, usage \Pic(X)
% \Picardstack_{X/B} denotes the Picard stack of X over B
% \Picardfunctor_{X/B} denotes the Picard functor of X over B
\def\Pic{\mathop{\mathrm{Pic}}\nolimits}
\def\Picardstack{\mathcal{P}\!\mathit{ic}}
\def\Picardfunctor{\mathrm{Pic}}
\def\Deformationcategory{\mathcal{D}\!\mathit{ef}}

% --- appended 2026-10-12 (Chief): the project own definitions of the reference/slogan/history
% display environments, copied mechanically from the shared Stacks preamble lines 8-12
% (\usepackage{verbatim} plus the three \newenvironment lines, which the chapter preamble keeps
% verbatim from the project preamble). Nothing added, removed or redefined: these environments
% carry no mathematical content and the project itself typesets them as comments. Without these
% three lines a counted statement whose own text uses them fails with Environment slogan undefined.
\usepackage{verbatim}
\newenvironment{reference}{\comment}{\endcomment}
\newenvironment{slogan}{\comment}{\endcomment}
\newenvironment{history}{\comment}{\endcomment}
\usepackage{cleveref}
\crefname{originallemma}{original lemma}{original lemmas}
\Crefname{originallemma}{Original lemma}{Original lemmas}
\makeatletter
\newlabel{algebra-lemma-bourbaki-fibres}{{05C2}{source}{}{sourcepointer-1}{}}
\newlabel{algebra-lemma-bourbaki-fibres@cref}{{[originallemma][0][]05C2}{[1][0][]source}{}{}{}}
\makeatother
\crefname{originalremark}{original remark}{original remarks}
\Crefname{originalremark}{Original remark}{Original remarks}
\makeatletter
\newlabel{algebra-remark-bourbaki}{{05E0}{source}{}{sourcepointer-2}{}}
\newlabel{algebra-remark-bourbaki@cref}{{[originalremark][0][]05E0}{[1][0][]source}{}{}{}}
\makeatother
\crefname{originallemma}{original lemma}{original lemmas}
\Crefname{originallemma}{Original lemma}{Original lemmas}
\makeatletter
\newlabel{lemma-universally-injective-to-completion}{{05DL}{source}{}{sourcepointer-3}{}}
\newlabel{lemma-universally-injective-to-completion@cref}{{[originallemma][0][]05DL}{[1][0][]source}{}{}{}}
\makeatother
\title{Universal injectivity of the map into the adic completion for a finite flat module, under a condition on the associated primes of the fibres}
\author{The Stacks Project Authors}
\date{stacks-project:flat.tex:local-snapshot-20261011; GNU FDL 1.2 or later}
\begin{document}
\maketitle
\noindent\textit{Editorial note.} All delivered mathematical text is the author's own, copied from the pinned Stacks Project chapter file flat.tex; no statement or proof step was rewritten, reordered or removed.\par\medskip
\noindent\textit{Editorial note.} Preamble provenance: the chapter's own retained minimal preamble (sources/stacks-flat/retained-preamble.tex, lines 1-85) is a verbatim excerpt of the native Xy-pic setup and of the native theorem declarations and mathematical macros, followed by the project's own reference/slogan/history display-environment definitions copied mechanically from the shared preamble; the wrapper class and the amsmath/amssymb/amsthm/graphicx base packages are supplied by the assembler, and extra\_packages is left empty because everything needed is already in the retained preamble. No macro was invented or redefined.\par\medskip
\noindent\textit{Editorial note.} Same-chapter pointers are located at their \textbackslash{}label lines inside the pinned chapter file; cross-chapter pointers are located in the project's own tag list staged in the snapshot (sources/stacks-chow/tags.txt), at the verbatim lines carrying their canonical tags.\par\medskip
\noindent\textit{Allowed standard prerequisites (editorial).}\begin{itemize}
\item Universal injectivity into the completion from a condition on associated primes of tensor products with field quotients. Same chapter flat.tex, label lemma-universally-injective-to-completion (invoked at proof line 2123). It is the general criterion, whose hypothesis is verified here in the restricted form stated.
\item The Bourbaki lemma and remark controlling the associated primes of tensor products by those of the fibres. Algebra, labels algebra-lemma-bourbaki-fibres and algebra-remark-bourbaki (invoked at proof lines 2126 and 2127). They guarantee that the associated primes of the tensor products \$N \textbackslash{}otimes\_R Q\$ are contained in the set of associated primes of the modules \$N \textbackslash{}otimes\_R \textbackslash{}kappa(\textbackslash{}mathfrak p)\$, so the hypothesis of the criterion is met.
\end{itemize}\medskip
\ifdefined\hypertarget\hypertarget{sourcepointer-1}{}\fi
\noindent\textit{Source reference.} Original lemma 05C2: sources/stacks-chow/tags.txt, source lines 6905--6905. Navigation pointer for \textbackslash{}ref\{algebra-lemma-bourbaki-fibres\} used inside the counted span (line 2126); the argument that consumes it is retained in the span.\par\medskip
\ifdefined\hypertarget\hypertarget{sourcepointer-2}{}\fi
\noindent\textit{Source reference.} Original remark 05E0: sources/stacks-chow/tags.txt, source lines 6973--6973. Navigation pointer for \textbackslash{}ref\{algebra-remark-bourbaki\} used inside the counted span (line 2127); the argument that consumes it is retained in the span.\par\medskip
\ifdefined\hypertarget\hypertarget{sourcepointer-3}{}\fi
\noindent\textit{Source reference.} Original lemma 05DL: sources/stacks-flat/pinned-original-source.tex, source lines 2060--2060. Navigation pointer for \textbackslash{}ref\{lemma-universally-injective-to-completion\} used inside the counted span (line 2123); the argument that consumes it is retained in the span.\par\medskip
\begin{lemma}
\label{lemma-universally-injective-to-completion-flat}
Let $R$ be a ring.
Let $I \subset R$ be an ideal.
Let $R \to S$ be a ring map, and $N$ an $S$-module.
Assume
\begin{enumerate}
\item $R$ is a Noetherian ring,
\item $S$ is a Noetherian ring,
\item $N$ is a finite $S$-module,
\item $N$ is flat over $R$, and
\item for any prime $\mathfrak q \subset S$ which is an associated prime of
$N \otimes_R \kappa(\mathfrak p)$ where $\mathfrak p = R \cap \mathfrak q$
we have $IS + \mathfrak q \not = S$.
\end{enumerate}
Then the map $N \to N^\wedge$ of $N$ into the $I$-adic completion of $N$
is universally injective as a map of $R$-modules.
\end{lemma}

\begin{proof}
This follows from
Lemma \ref{lemma-universally-injective-to-completion}
because
Algebra, Lemma
\ref{algebra-lemma-bourbaki-fibres} and
Remark \ref{algebra-remark-bourbaki}
guarantee that the set of associated primes of tensor products
$N \otimes_R Q$ are contained in the set of associated primes of
the modules $N \otimes_R \kappa(\mathfrak p)$.
\end{proof}

\end{document}
